% 03_assumptions.tex —— 模型假设
\section{模型假设}

为使模型可计算且与题意一致，本文作如下假设：
\begin{enumerate}
  \item 每台设备初始位置为其所属班组驻地，首次进入某车间作业前必须满足从班组驻地到该车间的转运时间。
  \item 同一台设备在同一车间内连续执行不同工序时，设备转运时间忽略。
  \item 同一台设备跨车间作业时，转运时间为车间间距离除以 $2\,\mathrm{m/s}$ 后向上取整。
  \item 对需要两类设备的工序，两类设备同时投入并分别完成该工序对应工程量；工序完成时间取两类设备完成时间的较晚者。
  \item 所有设备子任务的工作时间均按秒计，并由工程量和设备效率计算后向上取整。
  \item 问题四中，新增购置设备与同类型原有设备具有相同效率、移动速度和使用规则，且购买后立即可用于调度。
  \item 优化目标首先为最小化总完工时间；在总完工时间相同情况下，问题四进一步选择采购费用最小的方案。
\end{enumerate}

其中假设 4 对双设备协同工序的处理符合题目``不考虑两类设备先后顺序、不考虑两类设备之间等待时间''的表述。并且在问题三、四中，调度方案已经达到与该假设无关的 C 车间关键链下界，因此该处理不会影响最终最短工期结论。
